\documentclass[10pt]{article}
\usepackage[a4paper,margin=22mm]{geometry}
\usepackage{amsmath,amssymb,booktabs}
\usepackage[hidelinks]{hyperref}
\usepackage{url}
\usepackage{xurl}
\title{Technical Supplement: Activation-Aligned Asynchronous Residual MPC\\for Position-Controlled Manipulators with Learned Dynamics}
\author{Public evidence bundle for ROBIO 2026}

\begin{document}
\maketitle
\setlength{\textfloatsep}{8pt plus 2pt minus 2pt}
\setlength{\intextsep}{8pt plus 2pt minus 2pt}

\section{Purpose and evidence boundary}
This supplement provides implementation, configuration, and diagnostic detail
compressed in the conference manuscript. It introduces no results beyond the
frozen or deterministically reanalysed public evidence bundle.

The simulation results use MuJoCo position-controlled robot models. ABB
IRB~2400 is the primary robot; UR5e is an independently trained and calibrated
nominal within-robot replication. The bundle also contains a separate SO101
physical case study with a hardware-specific model and controller calibration.
The evidence does not establish zero-shot sim-to-real transfer,
hardware-independent generalization, external Cartesian accuracy, ABB
hardware-torque feasibility, hard-real-time scheduling, robust-MPC guarantees,
shared-model transfer, or arbitrary-morphology generalization. The post-freeze
effort sweep is diagnostic only.

\section{Controller and execution configuration}
The controller is a one-layer GRU (hidden size 256; history length 16) that
predicts velocity increments from aligned $[q_t,\dot q_t,q_t^{\rm ref}]$ tokens,
using one-step increment plus 20-step rollout loss on a 90/10 episode-disjoint
train/validation split; recurrent windows remain within their assigned episode
groups.
Frozen controller cost weights/scales are shared; XML models, data distributions,
checkpoints, normalizers, references, RobotSpec contracts, and delays are
robot-specific. The datasets are independently collected offline closed-loop
position-actuator transitions driven by non-MPC reference excitations: hold,
step, smooth-random, sinusoidal, delta-reference-random, and temporally
correlated random references. The latter uses an autoregressive random-reference
generator; no training rollout is generated by CEM or MPC. Workspace bounds,
normalized excitation scales, and optimizer settings remain robot-specific. For UR5e, the frozen dataset contains 120,000
episodes of 1,000 control steps (120 million transitions) collected at
normalized excitation scales 0.3, 0.5, 0.8, and 0.9. The complete
machine-readable contracts are released with the evidence bundle.

\begin{table}[h]
\centering\small
\caption{Frozen ABB controller settings. Per-joint vectors and all robot-specific
artifacts are in the manifests.}
\begin{tabular}{ll@{\qquad}ll}
\toprule
Setting & Value & Setting & Value \\
\midrule
Control interval & 10 ms & History length & 16 \\
Horizon / samples / iterations & 20 / 128 / 2 & Elite ratio & 0.08 \\
Replan interval / delay $D$ & 5 / 6 steps & Guard & 5 ms \\
Temporal discount & 0.95 & Projection & two-stage compiled \\
$w_q/w_{\dot q}/w_r/w_s$ & 1 / .10 / .20 / .05 &
$w_{\dot r}/w_{\ddot r}/w_{\rm first}$ & .05 / .02 / .20 \\
Joint / velocity barrier & 10 / 5 & Barrier max multiplier & 2 \\
Task position/orientation & 1 / .25 & Task scales & .05 m / .0873 rad \\
Feedback $K_q/K_{\dot q}/\delta r_{\max}$ & .3 / .015 / .015 rad &
Residual maximum & (0.12,.10,.12,.15,.15,.20) rad \\
\bottomrule
\end{tabular}
\end{table}

\section{Physical SO101 case study}
The physical result is a deployment validation of the residual-planning
principle, not a direct reuse of the six-DOF simulation hyperparameters. The
SO101 follower controls five arm joints at 30~Hz. Its separately collected GRU
uses 16 history tokens $[q,\dot q,q^{\rm ref}-q]$ and predicts a direct
$\Delta$-state target. The frozen hardware instantiation uses $H=6$, 128 CEM
samples, two iterations, residual authority $|r_j|\leq1^\circ$, and the
tracking-dominant objective $J=C_q$. The executable-command projector,
joint/velocity/acceleration/braking constraints, encoder quantization, and
startup/homing gates remain active. The formal matrix contains 54 held-out
trials and nine development-circle trials, with Direct IK, fixed Preview6, and
NN-MPC in rotated pre-registered order. Direct IK denotes direct execution of
the frozen IK-derived joint reference; Preview6 advances that same nominal by
six control steps. These predictor parameterizations are hardware-specific,
while residual CEM and asynchronous packet semantics are unchanged. All
technically complete trials are retained.

The primary endpoint is encoder-space joint RMSE over the frozen
\texttt{SEGMENT\_SHAPE\_LOOP} window. TCP values are forward-kinematics calculations
from encoder configurations using the calibrated fine MuJoCo model; no
external Cartesian tracker was used. The public aggregate and sanitized trial
ledger are available under \texttt{so101/analysis}.

\begin{table}[h]
\centering\scriptsize
\caption{Held-out SO101 physical aggregate. TCP is FK-derived and secondary;
command activity uses transmitted executable commands.}
\setlength{\tabcolsep}{3.0pt}
\begin{tabular}{lcccc}
\toprule
Method & Joint RMSE & FK TCP RMSE & Cmd. vel. RMS & Cmd. acc. RMS \\
 & (deg) & (mm) & (rad/s) & (rad/s$^2$) \\
\midrule
Direct IK & 0.910 & 8.912 & 0.0435 & 0.7699 \\
Fixed Preview6 & 0.675 & 6.538 & 0.0435 & 0.7701 \\
NN-MPC $\pm1^\circ$ & \textbf{0.353} & \textbf{3.715} & 0.0807 & 1.1985 \\
\bottomrule
\end{tabular}
\end{table}

Across the 54 held-out trials, NN-MPC reduced joint RMSE by 61.2\% versus
Direct IK and 47.7\% versus Preview6. It won all 18 matched held-out pairs
against both baselines; the paired means and confidence intervals are recorded
in \texttt{so101/analysis/public\_summary.json}. Across 20,956 physical planner
events, mean/p95/p99 latency was 22.73/28.36/29.93~ms (maximum 68.09~ms), with
three late-dropped solutions and no control deadline miss, planner failure, or
runtime safety violation. This timing result establishes deployability and
does not independently isolate activation-alignment necessity on hardware.
Diagnostic acceleration exceedances induced by encoder-command quantization are
reported separately in the trial ledger and are not executable-projector safety
violations.
Command velocity and acceleration RMS increased by 85.4\% and 55.7\% relative
to Direct IK, respectively, reproducing the simulation tracking--activity
trade-off.

\paragraph{Injected-delay stress.}
We additionally evaluated the same \texttt{ThreadedAsync} execution path with a
deterministic 33-ms delay inserted after packet post-processing and immediately
before publication. Controller, model, reference, CEM, residual, projection,
and safety settings remained frozen. An independent shadow pass produced 2,327
finite latency samples and selected $D=4$ from the measured p99.5 latency
($100.601$~ms) plus a 5-ms guard at 30~Hz; the resulting late-drop and
packet-expiry rates were both zero. The v2 OOD envelope covered 100.000\% of
executed histories and 99.683\% of both selected actions and predicted states.
The preceding v1 calibration was excluded because it had 1,734 samples and a
2.768\% late-drop rate.

\begin{table}[h]
\centering\scriptsize
\caption{Paired SO101 injected-delay stress results. Each row is a held-out
fast-ellipse phase block; the final row is the arithmetic mean across blocks.}
\setlength{\tabcolsep}{4pt}
\begin{tabular}{lccc}
\toprule
Block & Baseline & Delay33 & Difference \\
 & joint RMSE (deg) & joint RMSE (deg) & Delay33--baseline (deg) \\
\midrule
ellipse\_fast\_p0 & 0.34685 & 0.41896 & +0.07211 \\
ellipse\_fast\_p1 & 0.31772 & 0.45074 & +0.13302 \\
ellipse\_fast\_p2 & 0.33127 & 0.42876 & +0.09749 \\
Mean & 0.33195 & 0.43282 & +0.10087 \\
\bottomrule
\end{tabular}
\end{table}

Across the three matched blocks, planner-event timing statistics were computed
after concatenating all events within each condition. Baseline pooled latency
was 22.56~ms (mean), 28.33~ms (p95), and 29.37~ms (p99); under delay it was
79.32, 94.34, and 99.44~ms, respectively. The stressed p95 spanned 2.83
control periods. All six runs completed without control-loop deadline
misses, planner failures, late drops, packet expirations, command-limit
violations, or runtime safety violations. The delay condition was worse in all
three pairs; the paired Student-$t$ 95\% interval for the mean increase was
$[+0.02488,+0.17687]^\circ$ (three pairs, df=2). Encoder-quantization
acceleration diagnostics are reported separately and are not executable-
projector safety violations. This stress test supports physical deployability
and latency sensitivity of the aligned execution path, but it is not a hardware
\texttt{NaiveDelayed} or no-alignment ablation.

\paragraph{True-history window validation.}
Held-out model errors use fixed windows beginning at steps 16, 36, \ldots, 176
of each saved 200-step rollout. Every window therefore has a complete
16-token ground-truth recurrent context before its first predicted successor.
Each horizon is evaluated as a separate segmentwise RMSE: $h=1$ measures the
first successor and $h=20$ pools prediction steps $1{:}20$ within each window.
These are not terminal errors. The historical start-of-rollout validation,
which padded the missing initial GRU history, is retained in the evidence
bundle for audit but is not used in the manuscript table.

Residual packets encode corrections, not reusable absolute commands. At packet
age $\kappa=t-t_a$, the full protocol reconstructs and projects
\[
q_t^{\rm cmd}=\Pi_{v,\mathrm{acc},\mathcal Q}(\bar q_t^{\rm IK}+r_\kappa+\delta r_t;
q_{t-1}^{\rm cmd},\dot q_{t-1}^{\rm cmd}).
\]
The launch-anchor replay counterfactual instead projects the cached planned
absolute command $q_\kappa^{\rm plan}$. Its large failure is therefore evidence
for execution-time nominal-frame reanchoring required by an independently
evolving nominal, not a competitive controller comparison.

\paragraph{Optional post-selection model-disagreement gate.}
The released controller exposes an optional gate that is disabled
($\texttt{uncertainty\_mode=off}$) in every reported configuration. The planning
GRU alone generates and scores all CEM candidates. Only after selection,
independently trained, architecture-matched GRU dynamics models roll out the
same selected $H_u$-step joint-reference sequence; all use the same non-MPC
reference-excitation training protocol. If $\hat{x}^{(j)}_{h,d}$ is
the prediction of model $j$ at step $h$ and state coordinate $d$, and $s_d$ is
the planning model's training-state standard deviation, the gate records
\[
 u=\sqrt{\frac{1}{H_u n_x}\sum_{h=1}^{H_u}\sum_{d=1}^{n_x}
 \left(\frac{\operatorname{std}_j\hat{x}^{(j)}_{h,d}}{s_d}\right)^2}.
\]
This post-selection calculation neither rescales CEM candidates nor changes
their costs, ranks, or packet activation step. In monitoring mode it is
telemetry only. When the optional gate is enabled, a hysteretic rule can
attenuate residual authority after sustained disagreement. An incomplete
time-bounded evaluation, or high disagreement together with a predicted
joint-margin, residual-saturation, or tracking-error-growth condition, invokes
a strict fallback: both residual and packet feedback are set to zero, yielding
the same Projected Direct IK command. A limited drift state can instead use
current-reference feedback rather than packet-predicted feedback. Thresholds,
model count, short horizon, and evaluation deadline are configurable
implementation parameters; this paper does not calibrate or evaluate them and
claims no OOD-detection capability, closed-loop benefit, or safety guarantee.

\section{Artifacts, references, and calibration}
The manifests record robot contracts, command limits, model/data identities,
references, and activation delays. The public calibration payloads in
\url{abb/calibration} and \url{ur5e/calibration} use
\[
D=\left\lceil\frac{L_{\mathrm{P95}}+g}{\Delta t}\right\rceil,
\]
where $L_{\mathrm{P95}}$ is snapshot-to-publication latency, $g=5$~ms, and
$\Delta t=10$~ms. Independent 500-plan calibrations gave ABB
$L_{\mathrm{P95}}=52.28$~ms and $D=6$ control steps, and UR5e
$L_{\mathrm{P95}}=56.94$~ms and $D=7$ control steps. For UR5e,
\url{ur5e/contracts} additionally exposes the robot and reference contracts.

\section{Full diagnostics}
\begin{table}[h]
\centering\scriptsize
\caption{Held-out open-loop model errors from complete-history windows, pooled
over prediction steps $1{:}h$. $\sigma_u$ is normalized reference-excitation
standard deviation. Errors are $q/\dot q$ RMSE in rad/rad\,s$^{-1}$; Div. is
the divergent-window rate at $\sigma_u=0.8$.}
\setlength{\tabcolsep}{4pt}
\begin{tabular}{lccccc}
\toprule
Robot, $h$ & $q$ RMSE, $\sigma_u=.5$ & $\dot q$ RMSE, $\sigma_u=.5$ & $q$ RMSE, $\sigma_u=.8$ & $\dot q$ RMSE, $\sigma_u=.8$ & Div., $\sigma_u=.8$ \\
\midrule
ABB, 1 & .000670 & .001990 & .001069 & .005570 & 0\% \\
ABB, 20 & .004852 & .034472 & .006876 & .051398 & 0\% \\
UR5e, 1 & .000132 & .000465 & .000160 & .000532 & 0\% \\
UR5e, 20 & .001334 & .002745 & .001436 & .003023 & 0\% \\
\bottomrule
\end{tabular}
\end{table}

\begin{table}[h]
\centering\scriptsize
\caption{ABB retained-candidate replay diagnostics. Spearman is computed
within snapshots with defined ranks; all intervals are hierarchical bootstrap
95\% intervals.}
\setlength{\tabcolsep}{3pt}
\begin{tabular}{lcc}
\toprule
Metric & Estimate & 95\% interval \\
\midrule
Within-snapshot Spearman & 0.976 & [0.917, 1.000] \\
Pairwise concordance & 0.988 & [0.957, 1.000] \\
Relative realized regret & 0.192\% & [0, 0.481]\% \\
20-step branch $q$ RMSE & 0.889 mrad & [0.705, 1.121] mrad \\
\bottomrule
\end{tabular}
\end{table}

\paragraph{Model validation and CEM ranking.}
Tables~2--3 provide the principal held-out and replay diagnostics. Complete
window, rollout, excitation, and counterfactual-replay files are released in
the evidence bundle. Replay covers retained, projection-consistent online
candidates only; its lower branch error is not directly comparable with the
broader random-excitation validation.

\paragraph{Supporting planner diagnostics.}
Task-space reranking reduced deployed ABB TCP RMSE by 3.05~mm (95\% CI
[0.96, 4.32]~mm) relative to joint-only final-pool scoring, while increasing
end-to-end p95 latency by 6.82~ms. Two-stage compiled projection met its
pre-specified 5\% non-inferiority criterion against full compiled projection.
These are accuracy--latency diagnostics for supporting execution mechanisms,
not independent core tracking claims.

\paragraph{Robustness and timing.}
ABB robustness applies L3/L6 single-factor payload, actuator-gain, TCP-force,
and observation-noise perturbations. Per-condition/bootstrap results, timing,
packet behavior, projection activity, and paired Projected/Preview-IK
comparisons are released with the evidence bundle. Primary experiments had no
planner failure or imposed joint, command-velocity, or command-acceleration
violation; this does not imply a physical actuator limit.

\paragraph{Mechanism isolation.}
The frozen 20-case ABB matrix gives $+2.82$ mm for removing activation
alignment (95\% CI $[1.27,4.39]$ mm), $+654.60$ mm for launch-anchor absolute
replay ($[528.68,778.68]$ mm), and $-0.95$ mm for removing fast feedback
($[-2.25,0.28]$ mm), all relative to FullVirtual. A focused recurrent-history
isolation changed TCP RMSE by $-0.67$ mm with an interval crossing zero. Thus
the evidence supports activation consistency but does not establish an
independent nominal tracking gain from history propagation. Supporting
projection/task-score and post-freeze effort diagnostics did not replace the
frozen primary configuration.

\end{document}
